\section{Reflection and the duality group}
\label{rf:section}

We now close the chamber without losing its labelled middle-dimensional
classes, using the reflection framework of Davis
\cite[Sections~13 and~15, especially Remark~15.9]{Davis1983}.
A second construction will give the resulting group a proper
cocompact action on a CAT(0) space.  These are different developments:
the first retains the four-dimensional chambers, whereas the second is
used only for the assembly theorems.

The input from Proposition~\ref{ch:chamber} is a compact connected oriented
smooth four-manifold $D$, with connected boundary $S$ and free fundamental
group $H$, together with a map of pairs
\[
 (D,S)\longrightarrow(P,S).
\]
Here $D$ has the homotopy type of a finite graph with finitely many
two-spheres, its $\Z H$-intersection form has a specified framed
hyperbolic basis of $k$ planes, and $P$ is obtained by attaching
three-cells along all $2k$ basis spheres.  The pair $(P,S)$ is an
oriented Poincar\'e pair of dimension four and $P\simeq BH$.
The map on fundamental groups is the given identification with $H$.
We use a collar on $S$ in both spaces, inserting a mapping cylinder
on the $P$ side if necessary.

By Corollary~\ref{ch:boundary-homology}, the boundary map
$\pi_1(S)\to H$ is surjective.  In particular its associated cover of
$S$ is connected.

\subsection{Panels and finite developments}
\label{rf:panels}

Choose a finite flag triangulation $\mathcal L$ of $S$, for example a
barycentric subdivision of a smooth triangulation.  Write $V$ for its
vertex set and $S_s$ for the closed dual block of $s\in V$ in the
barycentric subdivision.  These blocks are three-dimensional panels.
For nonempty $T\subseteq V$, the intersection $S_T=\bigcap_{s\in T}S_s$ is nonempty exactly when
$T$ spans a simplex of $\mathcal L$; in that case it is a ball of
dimension $4-|T|$.  Collar neighborhoods of these intersections have
the usual orthant structure.

Let
\[
 C=\left\langle s\in V\ \middle|\ s^2=1,\quad
 st=ts\text{ whenever }\{s,t\}\in\mathcal L\right\rangle
\]
be the right-angled Coxeter group.  For a collared space $B$ with these
panels in its boundary, its development is
\[
 \begin{gathered}
 \mathcal U(C,B)=(C\times B)/\sim,
 \qquad T(x)=\{s:x\in S_s\},\\
 (c,x)\sim(c',x)\quad\Longleftrightarrow\quad
 c^{-1}c'\in C_{T(x)}.
 \end{gathered}
\]
Here $C_T$ is the subgroup generated by $T$ and $T(x)=\varnothing$
off the boundary.  Each chamber embeds: the equivalence relation never
identifies different values of $x$.  For nonempty $T(x)$ the subgroup
$C_{T(x)}$ is $(\Z/2)^{|T(x)|}$, and its sectors fill the corresponding
orthant neighborhood.

We use the elementary Coxeter word facts that special subgroups embed,
a finite special coset has a unique shortest element, and the right
descent set
\[
 \operatorname{Des}(c)=\{s\in V:|cs|<|c|\}
\]
is a clique.  These follow from the right-angled reduced-word rule of
commuting adjacent commuting letters and cancelling equal consecutive
letters; see also \cite[Chapters~4--5]{Davis}.

\begin{lemma}[Labelled finite developments]
\label{rf:shelling}
Put $Y=\mathcal U(C,D)$ and let $D_c$ be its chamber with label $c$.
Order $C$ by nondecreasing word length, breaking ties arbitrarily.
Every finite initial segment containing the identity is, after rounding
its collar corners, an iterated boundary connected sum of its chambers.
In particular
\[
 U_L=\bigcup_{|c|\leq L}D_c
\]
is a compact manifold with boundary the connected sum of the corresponding
copies of $S$.  Its fundamental group is the free product of their
chamber groups.  The same statements hold if any chambers have been
stabilized in their interiors.
\end{lemma}

\begin{proof}
The part of a new chamber $D_c$ already present is exactly
\begin{equation}
 \label{rf:incoming}
 A_c=\bigcup_{s\in\operatorname{Des}(c)}S_s.
\end{equation}
To verify the possible equal-length overlaps, consider a point with
panel set $T$.  Its incident chambers are the special coset $cC_T$.
Write $a$ for the shortest element of that coset.  Its members have
the form $a\prod_{s\in J}s$, with $J\subseteq T$, and length
$|a|+|J|$.  If a distinct incident chamber occurs no later than $c$,
then $c$ is not the shortest member, so it has a descent in $T$.
This places the point in \eqref{rf:incoming}.  The converse follows
from the shorter adjacent chamber across each descent panel.

The set in \eqref{rf:incoming} is a three-ball.  More generally a union
of dual panels indexed by a nonempty simplex is a ball.  One proves
this by adding its panels one at a time.  A new panel is a ball, and
its intersection with the preceding panels is the corresponding union
of dual panels in its boundary link.  Induction on dimension makes
that intersection a ball of one lower dimension.  The dual-block
coordinates give the product collars needed to glue the two balls.
This argument also applies to panels in any simplicial cover of $S$.

It remains to check that $A_c$ is a ball in the \emph{boundary} of the
preceding union.  Contractibility alone would not establish the claim.
At a face incident to $q$ panels, label the $2^q$ local orthants by
subsets of those panels, starting at the shortest incident chamber.
The already present orthants form a Boolean downset $I$: every proper
subset of a present subset has smaller word length.  This remains true
with arbitrary choices within a length tie.

When $I$ is nonempty and proper, its frontier is a PL hypersurface.
For an explicit verification, write local coordinates as
$x=u+t(1,\ldots,1)$ with $\sum u_i=0$, and label an orthant by its
positive coordinates.  Its union is described by
\[
 t\leq f(u),\qquad
 f(u)=\max_{A\text{ maximal in }I}\ 
               \min_{i\notin A}(-u_i).
\]
The function $f$ is finite, continuous, and piecewise linear.  Hence
the frontier is locally a graph.  Along the incoming faces some sectors
are present and the new chamber sector is absent, so these faces lie
on this frontier.  The panel identifications are PL, and the embedded
incoming three-ball has locally flat boundary in this three-manifold
frontier.  Gluing the new chamber along it is therefore a boundary
connected sum, after rounding.  Induction proves the manifold and
boundary assertions.  Van Kampen gives the labelled free product,
and interior stabilizations do not affect any collar argument.
\end{proof}

\subsection{The universal development}
\label{rf:universal}

Let
\[
 \Delta=\ker\bigl(C\longrightarrow(\Z/2)^V\bigr),
\]
where each generator maps to its own coordinate vector, and define
\[
 N=\Delta\backslash\mathcal U(C,D),\qquad
 Q=\Delta\backslash\mathcal U(C,P).
\]
The parity map is injective on every finite special subgroup.
Thus $\Delta$ has trivial intersection with every point stabilizer
in either development.  The quotient has finitely many chambers.
For $D$ the orthant charts show that $N$ is a closed four-manifold.
Give chamber $c$ the orientation $(-1)^{|c|}$ times the chosen chamber
orientation.  Adjacent chamber orientations fit across a panel, and
$\Delta$ preserves them.  Thus $N$ is oriented.  The same convention
orients the reflected Poincar\'e pairs.

Let $\widehat S\to S$ be the connected cover determined by
$\pi_1(S)\to H$, with lifted triangulation $\widehat{\mathcal L}$.
This is a flag triangulation.  Indeed, the vertices of a lifted clique
project to distinct pairwise adjacent vertices downstairs.  They span
a simplex there; lifts of its edges, and of its triangular relations,
place all the vertices in one lift of that simplex.  Denote by
$\widehat C$ the right-angled Coxeter group on
$\widehat{\mathcal L}$.  The deck action of $H$ permutes its generators.
Forgetting the lift gives an $H$-invariant homomorphism
$q:\widehat C\to C$, and hence a surjection
\begin{equation}
 \label{rf:group}
 \begin{gathered}
 \rho:\widehat C\rtimes H\longrightarrow C,
 \qquad \rho(\widehat c,h)=q(\widehat c),\\
 G=\rho^{-1}(\Delta).
 \end{gathered}
\end{equation}

Let $\widetilde P$ be the universal cover of $P$, with boundary
preimage $\widehat S$ and its lifted panels.  On
$\widehat Q=\mathcal U(\widehat C,\widetilde P)$ the action is
\[
 (\widehat c,h)[\widehat d,x]
     =[\widehat c\,h(\widehat d),hx].
\]
The map $[\widehat d,x]\mapsto[q(\widehat d),\overline x]$ descends
to the ordinary development.  On a chamber interior it is the usual
covering map.  At a boundary point it is also a covering map: the
lifted panel set maps bijectively to its downstairs panel set, and
the corresponding finite special groups map isomorphically.  These
local charts identify $\widehat Q/G$ with $Q$.  The action of $G$ is
free: if an element fixes a point, its $H$-component fixes the chamber
coordinate and is therefore trivial, after conjugating the chamber
label; the remaining stabilizer is a lifted finite special group,
on which $\rho$ is injective and whose image misses $\Delta$.

The space $\widehat Q$ is contractible.  Filter its chambers by word
length in $\widehat C$.  The first chamber $\widetilde P$ is
contractible.  Every further chamber meets the shorter ones in the
nonempty ball of its descent panels, by the argument of
Lemma~\ref{rf:shelling}.  Its attachment is therefore a homotopy
equivalence.  In an infinite word-length layer, new interiors are
disjoint and their mutual intersections already lie in shorter layers.
All these attachments may be performed simultaneously.  The CW
direct limit is weakly contractible and hence contractible.  This
proves, in particular, that $Q$ is a finite $K(G,1)$.

\begin{proposition}[The reflected pair]
\label{rf:development}
The space $Q$ is a finite oriented Poincar\'e complex of dimension four.
The natural map $f:N\to Q$ has degree one and induces an isomorphism
on fundamental groups, identifying both with the group $G$ in
\eqref{rf:group}.  If $b$ is the total number of chamber hyperbolic
planes in $N$, then
\[
 \pi_2(N)\cong(\Z G)^{2b}
\]
has the specified orthogonal hyperbolic basis, and
$\widetilde H_i(\widetilde N;\Z)=0$ for $i\ne2$.
The group $G$ is finitely presented, is torsion-free, and satisfies
\begin{equation}
 \label{rf:pd}
 H^i(G;\Z G)=0\quad(i\ne4),\qquad
 H^4(G;\Z G)\cong\Z
\end{equation}
with trivial right $G$-action.
\end{proposition}

\begin{proof}
Use the same universal development with $\widetilde D$ in place of
$\widetilde P$.  The preceding covering argument is unchanged.
The first chamber is simply connected; attaching subsequent simply
connected chambers along balls shows that this development is the
universal cover $\widetilde N$.  Its reduced homology is the direct
sum of the chamber reduced homologies, each concentrated in degree two.
There are finitely many chamber orbits in $N$.  For each one, induction
from its chamber subgroup gives
\[
 \Z G\otimes_{\Z H}\pi_2(D)
       \cong (\Z G)^{2k}.
\]
The spheres lie in chamber interiors, so different chambers are
orthogonal.  Within a chamber the form is its given $\Z H$-form,
extended to $\Z G$.  Reversing one vector in a plane absorbs any
chamber orientation sign.  This gives the asserted basis and form.

The construction of $P$ attaches exactly the indicated three-cells
away from the boundary.  Consequently $Q$ is obtained from $N$ by
attaching three-cells on this entire middle basis.  It follows at
once that $f$ induces the displayed identification of fundamental
groups.  These attachments split off a nonsingular middle summand
in both absolute and relative duality: on homology they quotient
by the sphere summand and on cohomology they restrict to its
annihilator.  With $[Q]=f_*[N]$, naturality of cap product identifies
the surviving duality maps with those of $Q$.  The finite free chain
criterion therefore gives Poincar\'e duality for $Q$, just as for the
chamber pair.  Equivalently one may glue the chamber duality maps
across the manifold collar sectors, where the orientations cancel
in pairs.  The top fundamental class is unchanged, so $f$ has degree
one.

For clarity, the regular-coefficient calculation is particularly
short.  In the cohomology sequence for these three-cell attachments,
the only potentially new map is evaluation
\[
 H^2(N;\Z G)\longrightarrow
       \operatorname{Hom}_{\Z G}(\pi_2(N),\Z G).
\]
Under duality it is the adjoint of the nonsingular hyperbolic form,
and hence an isomorphism.  The middle cohomology thus disappears,
whereas $H^4(Q;\Z G)=H^4(N;\Z G)=\Z$; all other degrees vanish.
The right action on this last group is trivial because the orientation
chosen above is preserved by $G$.

Finally, finiteness and asphericity of $Q$ give a finite presentation
and a finite free resolution of $\Z$ from the chains of $\widehat Q$.
They also give finite integral cohomological dimension, which excludes
torsion: restriction to a nontrivial finite cyclic subgroup would give
a finite projective resolution for that group, contrary to its
nonvanishing cohomology in arbitrarily high degrees.  This proves
all the assertions, including \eqref{rf:pd}.
\end{proof}

\subsection{A proper cocompact CAT(0) model}
\label{rf:assembly}

The lifted nerve has infinitely many vertices when $H$ is infinite.
Its Davis complex alone need not give a cocompact action of the
semidirect product.  The following support construction addresses this.

\begin{lemma}
\label{rf:cat0}
The group $\widehat C\rtimes H$, and hence its finite-index subgroup
$G$, acts properly and cocompactly by isometries on a proper CAT(0)
cube complex of dimension at most five.
\end{lemma}

\begin{proof}
Choose a locally finite tree $T$ on which the free group $H$ acts
freely and cocompactly, without inversions.  For the trivial group
one may take a single vertex.  Assign to each lifted-nerve vertex
$s$ a point $a_s\in T^{(0)}$, equivariantly under $H$.  There are
finitely many vertex and edge orbits in the lifted nerve.  Thus the
distances $d_T(a_s,a_t)$ for adjacent $s,t$ have a common finite bound.
Choose an integer $r$ large enough that the closed radius-$r$ subtrees
$T_s$ about $a_s$ intersect for every such pair.

For a vertex $v$ and an edge $e$ of $T$, put
\[
 A_v=\{s:v\in T_s\},\qquad
 A_e=\{s:e\subset T_s\}.
\]
These sets have uniformly bounded finite size: there are finitely many
generator orbits, and a fixed-radius ball in the free cocompact tree
contains uniformly finitely many orbit points.  Let $W_v$ and $W_e$
be the special subgroups of $\widehat C$ generated by $A_v$ and $A_e$.
For $e=[v,w]$ we have $A_e=A_v\cap A_w$, so the edge maps are special
subgroup inclusions.

The graph of groups over the tree $T$ has colimit $\widehat C$.
Indeed, occurrences of a generator $s$ are identified precisely over
its connected support $T_s$.  Every defining commutation of
$\widehat C$ occurs in some vertex group, because the two supports
intersect at a vertex.  Conversely, each vertex group uses the
\emph{induced} subgraph of the lifted nerve, and so introduces no
additional relation.  The resulting presentation is exactly the
right-angled presentation of $\widehat C$.  All the edge maps are
injective, as are the maps of these special subgroups into
$\widehat C$.

Let $\Sigma_v$ and $\Sigma_e$ be the Davis cube complexes of these
finite-generator right-angled groups.  They are CAT(0), and the special
subcomplexes $\Sigma_e\subset\Sigma_v$ are convex in the standard
piecewise Euclidean cube metric.  The CAT(0) assertion is
\cite[Theorem~12.2.1(i)]{Davis}; special-subgroup metric convexity is
\cite[Appendix, Proposition~A.2]{Swiatkowski}.
Form a tree of spaces with vertex spaces
$\widehat C\times_{W_v}\Sigma_v$ and edge cylinders
$\widehat C\times_{W_e}(\Sigma_e\times[0,1])$.
Its underlying tree is the Bass--Serre tree, whose vertices and edges
are the cosets $\widehat cW_v$ and $\widehat cW_e$, with their
indicated labels in $T$; it is not simply the original support tree.
Glue each cylinder at its endpoints by the special-subcomplex maps.
Convex gluing gives a CAT(0) space over every finite subtree, preserving
the factors as convex isometric subspaces
\cite[Theorem~II.11.1 and Remark~II.11.2]{BH}.
One can apply this gluing first over finite subtrees and then pass
to the union; the local finiteness verified next ensures completeness.

There is no local-finiteness problem from potentially infinite
Bass--Serre valence.  In a fixed vertex Davis complex, for any one
incident support edge, translates of its special subcomplex have
vertex sets equal to the cosets of $W_e$ in $W_v$.  These cosets
partition the Cayley vertices.  A finite subcomplex therefore meets
only finitely many of these translates.  Each support vertex has
only finitely many incident edges.  Since each $\Sigma_v$ is locally
finite, the whole tree of cube complexes is locally finite.  Its
cubes have unit edge lengths and uniformly bounded dimension, so
the induced length metric is proper and complete.

Left multiplication gives the $\widehat C$-action.  The equivariance
of the supports extends this to $\widehat C\rtimes H$.  An element
stabilizing a vertex space projects to an element of $H$ fixing its
support vertex, and that element is trivial.  Within the vertex
space its point stabilizer is therefore a finite Coxeter stabilizer.
The same reasoning applies to edge interiors.  The cellular action
on this locally finite complex is proper.  There are finitely many
$H$-orbits of support vertices and edges, and each corresponding
Davis complex has compact quotient by its special group.  Hence
the full action is cocompact.  Each clique of the lifted nerve has
at most four vertices; the vertex spaces have dimension at most four
and the edge cylinders dimension at most five.  Finally $G$ has finite
index by \eqref{rf:group}, so its restricted action has all the same
properties.
\end{proof}

It follows that $G$ satisfies the $K$- and $L$-theoretic Farrell--Jones
conjectures with coefficients: use Bartels--L\"uck
\cite[Theorem~B(ii)]{BL} and Wegner
\cite[Theorems~1.1 and~3.4]{Wegner}.  In particular, because $G$ is
torsion-free, the Whitehead group, reduced $K_0(\Z G)$, and negative
$K$-groups vanish, and the simple $L$-theory assembly map is an
isomorphism in every degree; see \cite[Proposition~0.3(i)]{BL}.
These are the assembly consequences used below.  No four-dimensional
rigidity or disk-embedding conclusion is being inferred from CAT(0)
geometry.

\subsection{The gallery map}

We finish by recording the geometric map needed in the truncation
argument. The preceding CAT(0) model supplies assembly; the map here
instead records positions and directions in the ordinary chamber
development. Let $\Sigma_C$ denote the ordinary Davis cube complex of
$C$.  Its vertices are the chamber labels $C$, and the link of the
identity vertex is $\mathcal L$, whose realization is $S$.

\begin{lemma}
\label{rf:davis-map}
There is a $C$-equivariant continuous map
\[
 \eta:Y\longrightarrow\Sigma_C
\]
which sends the inner core of each chamber to its label, respects
the panels, and has uniformly bounded image diameter on each chamber.
On the boundary collar of the central chamber its radial directions
at the identity vertex give the dual-block identification with $S$.
After interior chamber stabilizations, the same construction gives a
continuous map with the same panel, core, and diameter properties;
it is equivariant under any subgroup preserving the stabilized pattern.
\end{lemma}

\begin{proof}
The Davis chamber is the cone on the barycentric subdivision of
$\mathcal L$, with its usual dual panels on the base.  Identify that
base panelwise with $S$.  Map a boundary collar of $D$ to this cone
by that identification on its outer boundary and by coning to the
apex on its inner boundary; map the remainder of $D$ to the apex.
The maps agree across reflected panels and hence define $\eta$.
Every chamber image lies in the corresponding compact Davis chamber,
whose diameter is independent of its label.  At the identity apex,
radial projection of the cone base to a sufficiently small metric
sphere gives exactly its link identification.  Stabilization changes
only the collapsed inner core, so does not affect the construction.
\end{proof}
